\documentclass{article}
\usepackage{hyperref}
\title{Minimal LaTeX Docker Image}
\author{Marc Wäckerlin}
\date{\today}
\begin{document}
\maketitle

This project provides a minimal, reproducible LaTeX environment in a Docker container, following the philosophy of \href{https://github.com/mwaeckerlin/very-base}{mwaeckerlin/very-base}, \href{https://github.com/mwaeckerlin/nodejs}{mwaeckerlin/nodejs}, and \href{https://github.com/mwaeckerlin/nginx}{mwaeckerlin/nginx}:

\begin{itemize}
  \item \textbf{Minimalism:} Only the required binaries and files are included in the final image.
  \item \textbf{Security:} Runs as non-root, no shell or package manager in the runtime image.
  \item \textbf{Multi-stage builds:} Build in a larger image, copy only the output and dependencies into a minimal runtime image.
  \item \textbf{Reproducibility:} Consistent, predictable builds and outputs.
\end{itemize}

\section*{Usage}
\begin{enumerate}
  \item Place your LaTeX source (e.g., \texttt{sample.tex}) in the \texttt{doc/} directory.
  \item Run \texttt{docker compose up} to build the image and compile the LaTeX file to PDF.
  \item The resulting \texttt{sample.pdf} will be available in the \texttt{doc/} directory.
\end{enumerate}

\section*{Example}
A sample document is provided in \texttt{doc/sample.tex}. On \texttt{docker compose up}, it will be compiled to \texttt{doc/sample.pdf}.

\vfill
\noindent\textbf{License:} MIT

\end{document} 